%% ──────────────────────────────────────────────
%%  پیوست الف – خلاصه‌ی مدیریتی
%%  فایل: appendices/app-a-executive-summary.tex
%% ──────────────────────────────────────────────
\chapter{خلاصه‌ی مدیریتی}
\label{app:executive-summary}

%% ── معرفی ──
\begin{tcolorbox}[mybox, title={درباره‌ی این پیوست}]
این پیوست خلاصه‌ای غیرفنی و قابل‌فهم برای مخاطب عمومی — از جمله سیاست‌گذاران، روزنامه‌نگاران، فعالان مدنی، و دانشجویان — از یافته‌های کلیدی ۱۴ فصل اصلی کتاب ارائه می‌دهد. ارجاعات علمی سنگین حذف شده و تمرکز بر پیام‌های اصلی و توصیه‌های عملی است.
\end{tcolorbox}

%% ================================================================
\section{کتاب درباره‌ی چیست؟}
\label{sec:app-a-about}
%% ================================================================

این کتاب یک پرسش ساده اما بنیادین را دنبال می‌کند: \textbf{وقتی ملتی از وضعیتش ناراضی است و نیروی خارجی وعده‌ی «نجات» می‌دهد، آیا باید به آن امید بست؟}

برای پاسخ به این پرسش، بیش از پنجاه مورد مداخله‌ی خارجی — از ایران باستان تا اوکراین ۲۰۲۲ — بررسی شد. یافته‌ی اصلی: \textbf{در اکثریت قاطع موارد، مداخله‌ی خارجی نه‌تنها مشکل را حل نکرده، بلکه مشکلات جدید و اغلب بدتری ایجاد نموده است.}

اما این یافته به معنای «هیچ کاری نکنید» نیست. جایگزین نجات‌طلبی، \concept{عاملیت جمعی} است: ملت‌ها باید خودشان عامل تغییر باشند — و از کمک خارجی فقط به‌عنوان \textit{ابزار} (نه \textit{جایگزین}) استفاده کنند.


%% ================================================================
\section{خلاصه‌ی بخش اول: مبانی نظری}
\label{sec:app-a-part1}
%% ================================================================

\subsection{فصل ۱: چارچوب مفهومی}

«نجات‌طلبی» طیف گسترده‌ای دارد: از آرزواندیشی خالصانه‌ی مردم تحت ستم تا فرصت‌طلبی حساب‌شده‌ی قدرت‌ها. این کتاب هر دو سر طیف را بررسی می‌کند.

\subsection{فصل ۲: حقوق بین‌الملل}

حقوق بین‌الملل میان دو اصل متناقض گیر کرده: \textbf{حاکمیت دولت‌ها} (هیچ‌کس حق دخالت در امور داخلی دیگران را ندارد) و \textbf{حقوق بشر} (همه مسئول حمایت از انسان‌ها هستند). دکترین «مسئولیت حمایت» (R2P) تلاشی بود برای آشتی دادن این دو — اما در عمل ابزار قدرت‌های بزرگ شد:
\begin{itemize}
\item در لیبی ۲۰۱۱ اعمال شد (چون قذافی ضعیف بود)
\item در سوریه اعمال نشد (چون روسیه وتو کرد)
\item در یمن اعمال نشد (چون عربستان متحد غرب بود)
\item \textbf{نتیجه:} حقوق بشر بیشتر «امتیاز» است تا «حق» — اعمالش به منافع قدرت‌ها بستگی دارد
\end{itemize}

\subsection{فصل ۳: مدل تحلیلی}

هر مداخله از شش زاویه تحلیل شد:

\begin{table}[htbp]
\centering
\caption{مدل شش‌وجهی تحلیل مداخله (خلاصه)}
\label{tab:app-a-hexagonal}
\begin{adjustbox}{max width=\linewidth}
\begin{tabularx}{\linewidth}{>{\raggedleft\arraybackslash}p{3cm} >{\raggedleft\arraybackslash}X}
\toprule
\rowcolor{PrimaryDeep}
\textcolor{white}{\textbf{بُعد}} & \textcolor{white}{\textbf{پرسش کلیدی}} \\
\midrule
ساختاری-سیاسی & رژیم هدف چه نوعی است؟ بحران از کجا آمده؟ \\
\rowcolor{NeutralLight}
حقوقی-هنجاری & آیا مجوز حقوقی وجود دارد؟ \\
اقتصادی & چه منابعی در میان است؟ چه کسی سود می‌برد؟ \\
\rowcolor{NeutralLight}
نظامی-امنیتی & توازن قوا چگونه است؟ \\
روان‌شناختی-فرهنگی & ذهنیت مردم هدف و مداخله‌گر چیست؟ \\
\rowcolor{NeutralLight}
ژئوپولیتیک & منافع قدرت‌های منطقه‌ای و جهانی چیست؟ \\
\bottomrule
\end{tabularx}
\end{adjustbox}
\end{table}

نتایج مداخله در طیفی از A (رهایی کامل) تا F (استعمار یا الحاق) رتبه‌بندی شد. اکثر مداخلات معاصر در محدوده‌ی C تا E قرار می‌گیرند.


%% ================================================================
\section{خلاصه‌ی بخش دوم: تجربه‌ی ایران}
\label{sec:app-a-part2}
%% ================================================================

\subsection{فصل ۴: ایران باستان}
حمله‌ی اسکندر مقدونی هخامنشیان را نابود کرد اما فرهنگ ایرانی در نهایت فاتحان را جذب کرد. درس: \textbf{فتح نظامی لزوماً فتح فرهنگی نیست.}

\subsection{فصل ۵: سده‌های میانه}
حمله‌ی اعراب (قرن هفتم)، ترکان، و مغول (قرن سیزدهم) هر بار ایران را ویران کرد اما هر بار ایرانیان با حفظ زبان، فرهنگ، و ادبیات خود احیا شدند. درس: \textbf{مقاومت فرهنگی بلندمدت مؤثرتر از مقاومت نظامی کوتاه‌مدت است.}

\subsection{فصل ۶: صفوی تا قاجار}
عهدنامه‌ی ترکمانچای (۱۸۲۸) قفقاز را از ایران جدا کرد. قاجارها حاکمیت را تکه‌تکه از دست دادند — نه در یک جنگ بزرگ، بلکه در فرآیندی تدریجی. درس: \textbf{از دست دادن حاکمیت اغلب تدریجی و نامحسوس است.}

\subsection{فصل ۷: ایران معاصر}

\begin{table}[htbp]
\centering
\caption{چهار نقطه‌ی عطف مداخله در ایران معاصر}
\label{tab:app-a-iran-modern}
\begin{adjustbox}{max width=\linewidth}
\begin{tabularx}{\linewidth}{>{\centering\arraybackslash}p{1.5cm} >{\raggedleft\arraybackslash}p{3.5cm} >{\raggedleft\arraybackslash}X}
\toprule
\rowcolor{PrimaryDeep}
\textcolor{white}{\textbf{سال}} & \textcolor{white}{\textbf{رویداد}} & \textcolor{white}{\textbf{درس}} \\
\midrule
۱۹۴۱ & اشغال ایران توسط متفقین & حاکمیت ایران زیر پا گذاشته شد بدون توجه به بی‌طرفی \\
\rowcolor{NeutralLight}
۱۹۴۶ & بحران آذربایجان & شوروی «جمهوری خودمختار» ایجاد و سپس رها کرد \\
۱۹۵۳ & کودتای ۲۸ مرداد & آمریکا و بریتانیا دولت منتخب مصدق را سرنگون کردند \\
\rowcolor{NeutralLight}
۱۹۷۹ & انقلاب اسلامی & واکنش بلندمدت به مداخلات خارجی — اما خود نیز با مشکلات جدید \\
\bottomrule
\end{tabularx}
\end{adjustbox}
\end{table}

درس اصلی: \textbf{مداخله‌ی خارجی چرخه‌ی واکنش-رادیکالیزاسیون ایجاد می‌کند.} کودتای ۵۳ زمینه‌ساز انقلاب ۵۷ شد.


%% ================================================================
\section{خلاصه‌ی بخش سوم: تجربه‌ی جهانی}
\label{sec:app-a-part3}
%% ================================================================

\subsection{فصل ۸: اروپا}
اروپا خود هم عامل مداخله بوده و هم موضوع آن. انقلاب شکوهمند انگلستان (۱۶۸۸) نادرترین نمونه‌ی مداخله‌ی «موفق» است — اما شرایط آن تکرارناپذیر بود (نخبگان خودشان مداخله را مدیریت کردند).

\subsection{فصل ۹: استعمار و پسااستعمار}
مرزهای مصنوعی استعماری (خط‌کش‌هایی که بر نقشه‌ی آفریقا کشیده شدند) ریشه‌ی بسیاری از بحران‌های امروز هستند. استقلال‌های دهه‌ی ۱۹۶۰ اغلب «ظاهری» بودند: پرچم عوض شد اما ساختار وابستگی ماند.

\subsection{فصل ۱۰: مداخلات معاصر — خلاصه‌ی کلیدی}

\begin{table}[htbp]
\centering
\caption{کارنامه‌ی مداخلات معاصر (خلاصه)}
\label{tab:app-a-contemporary}
\begin{adjustbox}{max width=\linewidth}
\begin{tabularx}{\linewidth}{>{\raggedleft\arraybackslash}p{2.5cm} >{\raggedleft\arraybackslash}p{3cm} >{\centering\arraybackslash}p{1.5cm} >{\raggedleft\arraybackslash}X}
\toprule
\rowcolor{PrimaryDeep}
\textcolor{white}{\textbf{مورد}} & \textcolor{white}{\textbf{وعده}} & \textcolor{white}{\textbf{نتیجه}} & \textcolor{white}{\textbf{واقعیت}} \\
\midrule
عراق ۲۰۰۳ & آزادی و دموکراسی & C/D & جنگ فرقه‌ای، داعش، ۵۰۰٬۰۰۰+ کشته \\
\rowcolor{NeutralLight}
افغانستان ۲۰۰۱ & سرنگونی تروریسم & E & بازگشت طالبان پس از ۲۰ سال و ۲ تریلیون دلار \\
لیبی ۲۰۱۱ & حفاظت از غیرنظامیان & E & هرج‌ومرج، بازار برده، دو دولت رقیب \\
\rowcolor{NeutralLight}
سومالی ۱۹۹۲ & کمک انسان‌دوستانه & E & شکست و خروج، هرج‌ومرج ادامه‌دار \\
رواندا ۱۹۹۴ & — (عدم مداخله) & نسل‌کشی & ۸۰۰٬۰۰۰ کشته در ۱۰۰ روز \\
\rowcolor{NeutralLight}
سوریه ۲۰۱۱– & — (مداخله‌ی نیابتی) & D/E & ۵۰۰٬۰۰۰+ کشته، ۱۲ میلیون آواره \\
یمن ۲۰۱۵– & بازگرداندن دولت مشروع & E & بدترین بحران انسانی جهان \\
\rowcolor{NeutralLight}
اوکراین ۲۰۲۲ & — (مقاومت + کمک) & ؟ & مقاومت ملی قوی — ادامه‌دار \\
\bottomrule
\end{tabularx}
\end{adjustbox}
\end{table}

\textbf{الگوهای مشترک شکست:}
\begin{enumerate}
\item فقدان برنامه‌ی «روز بعد» — سقوط رژیم آسان، بازسازی غیرممکن
\item شکاف بین شعار (حقوق بشر) و انگیزه (نفت، ژئوپولیتیک)
\item اثر بومرنگ: عراق → داعش؛ لیبی → بحران مهاجرت؛ افغانستان → بازگشت طالبان
\item انتخابی بودن: در لیبی مداخله شد، در سوریه نه — تفاوت: منافع قدرت‌ها
\end{enumerate}


%% ================================================================
\section{خلاصه‌ی بخش چهارم: چرا نجات‌طلبی تکرار می‌شود؟}
\label{sec:app-a-part4}
%% ================================================================

\subsection{فصل ۱۱: دلایل روان‌شناختی}

سه مکانیزم اصلی:

\begin{tcolorbox}[defbox, title={سه مکانیزم روان‌شناختی نجات‌طلبی — زبان ساده}]
\textbf{۱. ناامیدی آموخته‌شده:}
وقتی سال‌ها تلاش برای تغییر بی‌نتیجه بوده، مردم باور می‌کنند «ما نمی‌توانیم» — و به بیرون چشم می‌دوزند.

\textbf{۲. آرزواندیشی:}
مردم تحت فشار آنچه را \textit{می‌خواهند} واقعیت باشد با آنچه \textit{هست} اشتباه می‌گیرند. «آمریکا/اروپا/سازمان ملل ما را نجات خواهد داد» — بدون توجه به شواهد تاریخی مبنی بر عکس آن.

\textbf{۳. بازتعریف هویت:}
وقتی مردم رژیم حاکم را «بیگانه‌تر از بیگانه» ببینند، ممکن است نیروی خارجی را «خودی‌تر از خودی» تلقی کنند — و مداخله‌اش را خوشامد بگویند.
\end{tcolorbox}

\subsection{فصل ۱۲: دلایل اقتصادی-ساختاری}

\begin{itemize}
\item \textbf{نظام جهانی نابرابر:} کشورهای «پیرامونی» (عراق، لیبی، افغانستان) ساختاراً آسیب‌پذیر در برابر مداخله‌اند
\item \textbf{نفرین منابع:} نفت هم انگیزه‌ی مداخله ایجاد می‌کند و هم ساختار سیاسی را فاسد می‌سازد
\item \textbf{صنعت جنگ:} شرکت‌های اسلحه‌سازی سالانه صدها میلیارد دلار از مداخلات سود می‌برند
\item \textbf{«بازسازی» سودآور:} تخریب و بازسازی هر دو سودآورند — برای مداخله‌گر
\end{itemize}

\subsection{فصل ۱۳: دلایل فرهنگی-ایدئولوژیک}

مداخله تقریباً همیشه با یک «پروژه‌ی فرهنگی» همراه است:
\begin{itemize}
\item میسیونرهای مسیحی + استعمار (قرن ۱۶–۲۰)
\item کمونیسم شوروی + مداخله در اروپای شرقی و جهان سوم
\item «دموکراسی‌سازی» آمریکایی + تغییر رژیم (عراق، لیبی)
\item صدور انقلاب اسلامی ایران + حمایت از حزب‌الله و اسد
\item وهابیسم سعودی + رادیکالیزاسیون اسلامی
\end{itemize}

\textbf{الگوی مشترک:} ابتدا ذهن‌ها فتح می‌شوند، سپس سرزمین‌ها.


%% ================================================================
\section{نتیجه‌گیری و توصیه‌ها}
\label{sec:app-a-recommendations}
%% ================================================================

\subsection{مدل نهایی: سه‌ضلعی نجات‌طلبی}

نجات‌طلبی محصول تعامل سه عامل است:

\begin{enumerate}
\item \textbf{ساختار اقتصادی نابرابر} → آسیب‌پذیری مادی
\item \textbf{مکانیزم‌های روان‌شناختی} → آمادگی ذهنی
\item \textbf{پروژه‌های فرهنگی-ایدئولوژیک} → مشروعیت‌سازی
\end{enumerate}

هر سه لازم‌اند. بدون آسیب‌پذیری مادی، نجات‌طلبی پایه ندارد. بدون آمادگی ذهنی، مردم مداخله را نمی‌پذیرند. بدون مشروعیت‌سازی فرهنگی، مداخله «تجاوز» به نظر می‌رسد نه «نجات».

\subsection{توصیه‌های عملی}

\begin{tcolorbox}[successbox, title={توصیه‌ها برای سیاست‌گذاران و فعالان مدنی}]
\textbf{برای ملت‌های تحت فشار:}
\begin{enumerate}
\item تفاوت «کمک‌خواهی» و «نجات‌طلبی» را بشناسید: کمک خارجی مشروع است — واگذاری سرنوشت به بیگانه مشروع نیست
\item نهادهای مدنی مستقل بسازید — نه وابسته به دولت و نه وابسته به خارج
\item تاریخ مداخلات خارجی را بیاموزید — حافظه‌ی تاریخی بهترین واکسن نجات‌طلبی است
\item مقاومت غیرخشونت‌آمیز را بیاموزید — دو برابر مؤثرتر از مقاومت مسلحانه
\item سوگیری‌های شناختی‌تان را بشناسید — آرزواندیشی دشمن واقع‌بینی است
\end{enumerate}

\textbf{برای جامعه‌ی بین‌المللی:}
\begin{enumerate}
\item قبل از هر مداخله‌ی نظامی، ابزارهای غیرنظامی را جدی بگیرید
\item «روز بعد» مهم‌تر از «روز حمله» است — بدون برنامه‌ی بازسازی، مداخله نکنید
\item انتخابی بودن مداخلات، مشروعیت کل نظام حقوق بشر را زیر سؤال می‌برد
\item تحریم‌های جامع، سلاح کشتار جمعی هستند — تحریم‌های هدفمند اخلاقی‌ترند
\item از صنعت جنگ حساب‌رسی کنید — مجتمع نظامی-صنعتی انگیزه‌ی ادامه‌ی جنگ دارد
\end{enumerate}
\end{tcolorbox}

\subsection{نتیجه‌گیری در یک جمله}

\begin{center}
\large\textbf{\rl{آزادی‌ای که از بیرون بیاید، آزادی نیست.}}\\[6pt]
\large\textbf{\rl{و ملتی که خود را نجات نداده، نجات نیافته است.}}
\end{center}

\cleardoublepage